\section{Geschäftsmodell und Besonderheiten}
\label{sec:geschaeftsmodell}

Die Allianz SE mit Hauptsitz in München zählt zu den weltweit führenden
Versicherungsunternehmen und Vermögensverwaltern. Im Geschäftsjahr 2025
beschäftigte die Allianz Gruppe 156.875 Mitarbeitende und erzielte ein
Geschäftsvolumen von rund 186,9 Milliarden Euro
\citep{allianzAnnualReport2025}. Das operative Geschäft gliedert sich in die
Segmente Schaden- und Unfallversicherung, Lebens- und Krankenversicherung
sowie Asset Management. Zum letztgenannten Segment gehören unter anderem die
Vermögensverwalter PIMCO und Allianz Global Investors. Im Mittelpunkt dieser
Arbeit steht das B2C-Privatkundengeschäft im Versicherungsbereich.

\begin{table}[H]
  \centering
  \caption{Vereinfachte Gegenüberstellung zentraler Einnahmen und
    Kostenblöcke im B2C-Versicherungsgeschäft}
  \label{tab:b2c-geschaeftsmodell}
  \begin{tabularx}{\textwidth}{@{}XX@{}}
    \toprule
    Einnahmen & Kostenblöcke \\
    \midrule
    Beiträge und Prämien der Privatkundinnen und Privatkunden
      & Schadens- und Leistungsaufwendungen \\
    Kapitalanlageerträge
      & Personal-, Vertriebs-, Service- und Betriebskosten \\
    \bottomrule
  \end{tabularx}
\end{table}

Die Profitabilität des Versicherungsgeschäfts hängt wesentlich von einer
risikoadäquaten Prämienkalkulation ab. Die vereinnahmten Prämien müssen neben
den erwarteten Schadens- und Leistungsaufwendungen auch Betriebs- und
Kapitalkosten langfristig decken. Eine belastbare Risikoeinschätzung mithilfe
geeigneter Datenanalysen besitzt daher eine hohe betriebswirtschaftliche
Bedeutung.
